\chapter{Systems Analysis and Design}

\section{Systems Analysis}
% Analyze the operational context and system requirements.

\subsection{Requirements Analysis}
% Outline the specific requirements gathered for the project.

\subsubsection{Functional Requirements}
% List the features, behaviors, and functions the system must perform.

\subsubsection{Nonfunctional Requirements}
% List the performance constraints, security, and quality attributes of the system.

\subsection{Feasibility Analysis}
% Provide an overview of the feasibility assessment for your project.

\subsubsection{Technical Feasibility}
% Evaluate whether the required tools and technologies are attainable.

\subsubsection{Operational Feasibility}
% Assess how well the proposed system will meet user needs and operational workflows.

\subsubsection{Financial Feasibility}
% Analyze the cost, budget, and economic viability of your project.

\subsubsection{Schedule Feasibility}
% Evaluate the timeline and ensure the project can be completed within the allocated time frame.

\subsection{Data Modeling}
% Illustrate how data is structured, organized, and related within the system.

\subsection{Process Modeling}
% Describe the data flow, workflows, and business processes handled by the system.

\section{Systems Design}
% Detail the structural blueprint and technical architecture of the system.

% This subsection is required for PRJ 281 only.
\subsection{Architectural Design}
% Describe the highlevel software architecture, components, and their interactions.

% This subsection is required for PRJ 281 only.
\subsection{Database Schema Design}
% Detail the database tables, fields, relationships, and constraints.

\subsection{Interface Design}
% Illustrate the user interface layouts, navigation, and user experience design.

% This subsection is required for PRJ 281 only.
\subsection{Physical Model}
% Describe the physical deployment, hardware, and infrastructure setup.
